\section{Pseudo-injective estimation: replacing the flow head with a regularized VAE}
\label{sec:pseudoinj}

\Cref{sec:flow} proposed a flow head to get an exactly invertible $g$.
It was abandoned before implementation: an exactly invertible map cannot reduce dimension
(\cref{sec:flow:motivation}), so a square injective flow head is structurally the wrong shape for M1, 
which needs roughly ten output latents, not a $150{,}528$-dimensional one. This section proposes what 
replaces it: not an exact diffeomorphism, but a pair of networks -- a frozen, pretrained backbone and a 
trainable head -- each carrying an explicit, measurable regularizer that keeps them \emph{close to} 
\textbf{injective}, while actually reducing dimension the way \nameref{sec:archB} does. 
We first restate the problem this is meant to solve, then the alternatives that were considered and set 
aside, then the two pieces of the proposed pipeline (\cref{sec:pseudoinj:invae,sec:pseudoinj:blae}), 
then a first, wrong version of the
training objective and why it had to be corrected (\cref{sec:pseudoinj:jacobian}), then the corrected
architecture (\cref{sec:pseudoinj:arch}), and finally an explicit list of gotchas
(\cref{sec:pseudoinj:gotchas}) that a reader implementing this should not discover the hard way.

\subsection{The problem this section answers}
\label{sec:pseudoinj:problem}
\Cref{sec:archB:gap} already states the gap in \nameref{sec:archB} plainly: a frozen ResNet-50 is
provably non-injective (any pooling convnet is), so nothing constrains $\mathrm{BB}\circ g^{-1}$ to be
injective, differentiable, or otherwise well-behaved, and the gap is \emph{unexamined} -- nothing in the
pipeline measures how bad it is. \Cref{sec:flow} closed that gap exactly, at the cost of being unable to
reduce dimension. What is wanted is a middle position: a backbone and a head that are not exactly
injective, but whose deviation from injective is bounded and estimated, not ignored, while still landing
at M1's small latent dimension (\emph{``pseudo-injective''}).

\subsection{Alternatives considered and set aside}
\label{sec:pseudoinj:discarded}
\begin{itemize}[leftmargin=1.5em]
\item \textbf{Rectangular (injective) flows} \cite{caterini2021}. These are the flow family that can, in
      principle, reduce dimension while staying exactly injective, by computing the Gram determinant
      $\det(J^\top J)^{1/2}$ instead of a square Jacobian determinant (\cref{sec:flow:injective}). Set
      aside: the published results barely reach CIFAR-10 scale, not ImageNet-scale full-resolution 
      images, and computing (or differentiating) the Gram determinant is itself an open research 
      problem, not an engineering cost.
\item \textbf{A $(z,w)$ padding variant.} Keep an exactly bijective, dimension-preserving flow as in
      \cref{sec:flow}, but declare only the first $\sim\!10$ coordinates of its output ``$z$'' (M1's
      latent) and treat the rest as a nuisance padding vector $w$, discarded after training. Set aside:
      this sits outside \cref{thm:ng} as written (the theorem is stated for $Z$ itself, not a marginal or
      a projection of a larger vector with no constraint tying $w$ to anything), and it does not obviously
      buy anything \nameref{sec:archB} does not already give more cheaply.
\end{itemize}

\subsection{IN-VAE: a pretrained, frozen, reconstruction-trained backbone}
\label{sec:pseudoinj:invae}
\textbf{What it is.} A large, pretrained image autoencoder of the kind used as the first stage of latent
diffusion pipelines: an encoder/decoder pair trained end-to-end on a reconstruction (plus, depending on
the variant, adversarial and/or perceptual) objective at ImageNet scale, whose \emph{encoder} we reuse
frozen as BB, exactly as \nameref{sec:archB} reuses a frozen ResNet-50. We call this class \textbf{IN-VAE}
in this note: pseudo-injective not because any Jacobian argument certifies it, but because reconstruction
training directly rewards the encoder-decoder pair for approximately undoing each other, which a
classification objective (plain ResNet-50, or the supervised-pretraining case of \cref{sec:pretrain}) has
no reason to do at all. Unlike \nameref{sec:archB}'s pooled $\sim\!3072$-dim vector, IN-VAE's output is a
\emph{spatial} latent, a tensor of shape $C\times h\times w$ (concretely $32\times16\times16$ for the
chosen candidate, \cref{sec:pseudoinj:arch}), not a flattened, globally pooled embedding -- this is
relevant to the head architecture below.

\textbf{Candidates surveyed.} SD-VAE kl-f8 (embedding width $4096$),
\footnote{the regularization type on the latent: KL-regularized, $\times$8 spatially downsampled}
the IN-VAE/REPA-E baseline used in
this note ($8192$), E2E-VAE ($8192$-class, but its representation-alignment loss is trained to match a
downstream generative model's preferred feature geometry, which is exactly the kind of pressure
\cref{sec:pretrain} already flags as likely to squeeze out style/location cues), FLUX/SD3 VAE
($16384$), DC-AE, and MAETok/STELLAR (semantically-shaped tokenizers, same style-squeezing risk as
E2E-VAE, for a different reason: their objective explicitly rewards semantic/content compression). Every
candidate carrying this alignment or semantic-shaping risk -- E2E-VAE, MAETok, STELLAR -- must be screened
with the linear-probe diagnostic of \cref{app:probe} (species accuracy plus within-species location
accuracy) before it is used, exactly as \cref{sec:pretrain} already requires for the ResNet-50 case; none
of them is used or discussed further in this note.

\textbf{Why IN-VAE (the REPA-E-style baseline) \cite{leng2025repae} is the default here.} It is trained on plain reconstruction
(plus generic perceptual/adversarial terms), with no alignment or semantic-shaping loss pulling it toward
discarding style/location information; its embedding width ($8192$, i.e.\ $32\times16\times16$) is
comparable to the other non-flagged candidates; and, as an encoder-decoder pair with a released,
well-tested implementation (Hugging Face \texttt{diffusers.AutoencoderKL} \cite{aekl}), it needs no 
bespoke pipeline code, only a frozen forward pass -- the same iteration-speed argument \cref{sec:archB} 
already makes for a frozen backbone. Everything from here on assumes this choice; it should still pass 
\cref{app:probe}'s screen before being trusted, the same as any candidate.

\subsection{BLAE: an injective/bi-Lipschitz regularizer for the trainable head}
\label{sec:pseudoinj:blae}
\textbf{What it is.} Zhan et al.\ \cite{zhan2026blae} assume that high-dimensional data resides
on a $k$-low-dimensional manifold $\Mk \subset \R^m,\ k = \dim(\Mk)$ and that there's an autoencoder 
that learns this intrinsic representation through a pair of parametrized mappings 
$(\E_\theta,\D_\phi)$:
\begin{itemize}
\item The encoder $\E_\theta : \R^m \to \R^n (k < n \ll m)$ compresses data onto a low-dimensional 
      latent space;
\item The decoder $\D_\phi : \R^n \to \R^m$ reconstructs the original data from the latent 
representation.
\end{itemize}

They propose two losses to add to an otherwise ordinary autoencoder, meant to push its encoder 
toward being injective without requiring an exactly invertible architecture:
\begin{itemize}[leftmargin=1.5em]
\item an \emph{injective} (separation) loss, which penalizes distinct inputs landing too close 
      together in latent space (their Theorem~1 gives a separation criterion);
\item a \emph{bi-Lipschitz} relaxation term, which additionally penalizes local distortion of
      \emph{geodesic} distances between the ambient data (estimated once via an Isomap-style 
      shortest-path computation) and Euclidean distances between the corresponding latents; the 
      paper argues (their Theorems~3-4) this relaxation is admissible at $O(n)$ latent dimensions, 
      where a hard embedding guarantee would need $O(k^2)$.
\footnote{According to the Nash embedding theorem \cite{nash1956}, a $k$-dimensional compact 
Riemannian manifold requires a latent dimension of $O(k^2)$ to guarantee an isometric embedding
(distance preserving); the required latent dimension $n$ satisfies 
$\frac{k(k+1)}{2} \le n \le \frac{k(3k+11)}{2}$.}
\end{itemize}
This is a \emph{trained regularizer}, not a proof of injectivity, and it is unproven at ImageNet 
scale: the paper's own experiments are Swiss Roll, dSprites, MNIST, CIFAR-10, Bird, and Teapot.

\textbf{What is reusable from the released repository, and what is not.} The repository's dataset-
specific autoencoder classes are toy-sized and toyishly-implemented (latent dims $2$--$16$, small 
fixed image sizes) and not directly usable.
What \emph{is} architecture-agnostic and directly reusable: the injective loss and the bi-Lipschitz 
term themselves (they take only a latent batch and a precomputed geodesic-distance table, nothing
dataset-specific), and the geodesic-distance-table computation (\cref{sec:pseudoinj:gotchas} 
discusses its cost at TI scale). \textbf{No VAE variant is released}: a stochastic BL-VAE is 
described only in the paper's Appendix~C.6.2, demoed on Swiss Roll, and never shipped in the code 
(a repository search for VAE/reparameterization/log-variance/KL code returns nothing in Python 
files; there's a Jupyter notebook with some VAE-denoted variables, but their examination shows that 
this is not a true VAE).

\subsection{Why a VAE head, and not a deterministic autoencoder}
\label{sec:pseudoinj:jacobian}
\textbf{The question.} The exact change-of-variables objective of \cref{eq:flow-nll} needs a volume
term. Neither IN-VAE's encoder nor a BLAE-regularized head is a square, exactly invertible map, so the
term is the Gram determinant $\det(J^\top J)^{1/2}$ of \cref{sec:flow:injective}, not an ordinary
determinant. The question is whether this objective can be used for the head. An earlier version of this
note said that $J$ ``cannot be computed''. That statement is too strong, and the justification below
replaces it.

\textbf{What is and is not the obstacle.}
\begin{itemize}[leftmargin=1.5em]
\item \emph{Computing $J$ is not the obstacle.} The trainable decoder maps a latent of dimension $10$ to
      the image-side tensor, so its Jacobian has only $10$ columns and $J^\top J$ is a $10\times10$
      matrix. Forward-mode differentiation gives $J$ with about $10$ passes per sample. This is exactly
      what the bi-Lipschitz term of \cref{sec:pseudoinj:blae} already does when it evaluates the decoder
      Jacobian. That term is a regularizer, however, and only needs approximate singular values. A
      likelihood needs the exact $\log\det$ and its gradient at every step.
\item \emph{Cost is a real but secondary obstacle.} Backpropagating through $\log\det(J^\top J)$ means
      differentiating through forward-mode Jacobians of a convolutional decoder with attention, at every
      training step and for every sample. The attention blocks must also be run with the unfused (math)
      backend, because the fused kernels have no forward-mode rule. This is affordable for a
      regularizer weighted by a small coefficient, but it is a heavy price for the main objective.
\item \emph{The Gram determinant does not factorize over stages.} For square maps, $\log|\det|$ of a
      composition is the sum of the per-stage terms, which is what made the flow section tractable
      (\cref{eq:flow-jac-split}). For non-square maps it is not. The volume term of the composite
      ``frozen IN-VAE decoder after head decoder'' is not the sum of a backbone term and a head term, so
      the frozen backbone cannot be handled separately.
\item \emph{The main obstacle: the density is degenerate.} A deterministic decoder from a
      $10$-dimensional latent has an image that is a $10$-dimensional manifold inside a much
      larger space. The pushforward of $p_{\hat Z}$ is then supported on that manifold and has no density
      with respect to Lebesgue measure on the large space. The change-of-variables formula yields a
      density only with respect to the surface measure of the manifold, and it is meaningful only for
      points that lie on it. Real images (or their cached embeddings) never lie exactly on it, so
      $\log p_X(x)$ is undefined without an additional off-manifold noise model. This holds no matter how
      cheaply the Gram determinant could be computed.
\end{itemize}

\textbf{The rejected first candidate.} 
The first idea was to keep BLAE's deterministic autoencoder
and evaluate M1's prior at the single encoded point $\hat z=\mathrm{Enc}(x)$, as in
\cref{eq:flow-prior-term}, with no volume term. This silently replaces the target $\log p_X(x)$ by
$\log p_{\hat Z}(\hat z\mid u)$, which is a poor proxy to the extent that the encoder is non-injective.
Nothing in that loss bounds or even reports the gap.

\textbf{Decision: use a VAE for the trainable head.} A VAE needs neither the Gram determinant nor the
manifold assumption. The Gaussian decoder with fixed variance is exactly the off-manifold noise model
that makes the density on the large space well defined, and the ELBO then requires no Jacobian at all.
This is why Ng et al.'s estimator (\cref{app:code}) and \nameref{sec:archA}/\nameref{sec:archB} all use
an ELBO. The price is the usual, already accepted approximation: $q(\hat Z\mid X)$ only approximates the
true posterior, and the reconstruction term only penalizes the cases where decoding fails to invert
encoding. The Jacobian then appears in exactly one place, the BLAE bi-Lipschitz regularizer, where an
approximate, weighted use is sufficient.

\subsection{The resulting architecture}
\label{sec:pseudoinj:arch}
\begin{enumerate}[label=(\arabic*),leftmargin=2em]
\item \textbf{Frozen backbone.} IN-VAE's pretrained encoder (\cref{sec:pseudoinj:invae}), run once 
       per image, cached, exactly as \nameref{sec:archB} caches $\tilde x=\mathrm{BB}(x)$ -- except 
       here $\tilde x$ is the spatial IN-VAE latent, a $32\times16\times16$ tensor, not a pooled 
       ResNet vector.
%-----------------
\item \textbf{Trainable head.} A \emph{fresh, untrained} instance of Hugging Face's
      \texttt{diffusers.AutoencoderKL} \cite{aekl} (not a pretrained checkpoint), 
      used \emph{unmodified} -- rather 
      than editing its \texttt{quant\_conv} and \texttt{post\_quant\_conv} internals, two small 
      trainable wrapper blocks sit outside it and convert between its native spatial latent and M1's 
      flat vector.

      \textbf{Sizing the stock module.} \texttt{AutoencoderKL} \cite{aekl} is configured, but not 
      altered internally, to fit the $32\times16\times16$ IN-VAE input:
      \begin{itemize}[leftmargin=1.3em]
      \item In diffusers the last encoder stage does not downsample, so \texttt{N} stages give 
            \texttt{N-1} halvings.
            
            \texttt{down\_block\_types = ("DownEncoderBlock2D", "DownEncoderBlock2D",}\\ 
            \texttt{"DownEncoderBlock2D")} (3 stages,
            $16\to8\to4$), with the symmetric \texttt{up\_block\_types} triple;
      \item \texttt{block\_out\_channels = (32, 64, 128)};
      \item \texttt{layers\_per\_block = 1};
      \item \texttt{in\_channels = out\_channels = 32};
      \item \texttt{latent\_channels} left at whatever stock value keeps \texttt{quant\_conv}/
            \texttt{post\_quant\_conv} well-defined (e.g.\ $4$); this is \emph{not} M1's $\sim\!10$-
            dim latent -- that reduction happens entirely in the wrapper blocks below, outside the 
            module.
      \end{itemize}
      With this config, stock \texttt{AutoencoderKL.encode} returns its own spatial latent 
      distribution at $4\times4\times4$ (i.e.\ \texttt{latent\_channels}$=4$, spatial size 
      $4\times4$), unchanged from the library's normal behavior.

      \textbf{Pre-block (after \texttt{AutoencoderKL.encode}, before M1).} Flatten (or global-average-
      pool) the module's own $4\times4\times4$ output, then a single linear layer maps this flat vector 
      to M1's $20$ scalars $(\hat\mu,\hat{\log\sigma^2})$ (i.e.\ $10$ content/style latents). This 
      linear layer is the actual dimensionality-reducing step; \texttt{AutoencoderKL} itself is never 
      told about, and never computes, the $10$-dim target.

      \textbf{Post-block (after M1's prior/loss, before \texttt{AutoencoderKL.decode}).} A single linear
      layer expands the sampled $10$-dim $\hat z$ back to a flat $4\times4\times4$ vector, reshaped to
      $4\times4\times4$ and passed into stock \texttt{AutoencoderKL.decode} unmodified.

      \textbf{Why this is preferred over editing the module.} The pre/post linear layers are the only
      hand-written, unvalidated pieces; \texttt{AutoencoderKL}'s own encode/decode path, KL
      reparameterization, and internal conv blocks stay exactly as tested and released, with no risk of 
      a silent behavioral change from touching \texttt{quant\_conv}/\texttt{post\_quant\_conv} directly. 
      The cost is one extra pair of thin layers and one more bottleneck seam (module output $\to$ M1 
      vector) to account for when tracing where distortion enters the pipeline 
      (\cref{sec:pseudoinj:gotchas}).
%-----------------
\item \textbf{M1 prior, unchanged.} The content/style ANM prior of \cref{sec:use} -- \texttt{mean\_net},
      \texttt{log\_var\_net}, the adjacency \texttt{adj\_mat}, sink-freezing -- sits on top of the head's
      output exactly as in \cref{app:code}; nothing about the prior side changes here, only the recognition
      network $q(\hat Z\mid X)$ feeding it.
\item \textbf{BLAE losses, added to the head's training, computed on $\hat\mu$.} The injective (separation)
      loss and the bi-Lipschitz term of \cref{sec:pseudoinj:blae} are added to the head's loss, evaluated
      on the posterior mean $\hat\mu$ (never on a noisy sample $\hat z=\hat\mu+\epsilon$, since the
      geodesic-preservation objective is inherently a statement about a deterministic map -- the
      reparameterization noise must not enter this term). The geodesic-distance table (\texttt{GeoD}) is
      built once from the cached IN-VAE embeddings of the training-location images.
\end{enumerate}
The resulting per-image loss, replacing \texttt{Trainer.step} of \cref{app:code} on the head side:
\begin{equation}
\begin{aligned}
\mathcal L(x,u) &= \lambda_{\mathrm{kl}}\,\mathrm{KL}\bigl(q(\hat Z\mid x)\,\|\,p_{\hat Z}(\hat Z\mid u)\bigr) \\
&+ \lambda_{\mathrm{rec}}\,\mathrm{recon}(\tilde x,\hat{\tilde x}) \\
&+ \lambda_{\mathrm{inj}}\,\ell_{\mathrm{inj}}(\hat\mu) \\
&+ \lambda_{\mathrm{biLip}}\,\ell_{\mathrm{biLip}}(\hat\mu,\mathtt{GeoD}) \\
&+ \lambda_{\mathrm{sparse}}\,\ell_{\mathrm{sparse}}(\hat A) \\
&+ \lambda_{\mathrm{moral}}\,\ell_{\mathrm{moral}}(\hat A, A^{t-1}),\footnotemark\\
&\quad A^{t-1}\ \text{is the binary adjacency matrix estimated in the previous} \\
&\quad \text{(i.e.,}\ (t\!-\!1)\text{-th) iteration}.
\end{aligned}
\label{eq:pseudoinj-loss}
\end{equation}
\footnotetext{Ng implementation dropped the moral loss term in their implementation --
see \cref{app:code:lossterms}.}%
where $\tilde x$ is the cached IN-VAE embedding (the reconstruction target, one level removed from the raw
image, as in \nameref{sec:archB}), and the first, second and fifth terms are exactly \cref{app:code}'s
\texttt{loss\_kl}, \texttt{loss\_rec} and \texttt{loss\_sparsity}. \emph{No Jacobian term appears anywhere
in \cref{eq:pseudoinj-loss}}, by the decision of \cref{sec:pseudoinj:jacobian}.

\subsection{Why this is more principled than plain Option B -- and what it still is not}
\label{sec:pseudoinj:principled}
Both stages of this pipeline carry an explicit, measurable regularizer toward injectivity: IN-VAE's
reconstruction training (rewards approximate invertibility directly, unlike a classification objective);
the head's injective and bi-Lipschitz losses (bound local distortion of distances, with a stated theoretical
target). Plain \nameref{sec:archB} has neither -- \cref{sec:archB:gap}'s gap is totally unexamined there.
This is genuinely better: instead of an unknown, unbounded gap, the gap here is a quantity the bi-Lipschitz
constants put a number on.

It is not a proof of injectivity, and \cref{thm:ng}'s hypothesis -- $g$ an exact $C^2$ diffeomorphism onto
its image -- is still, strictly, violated by both IN-VAE and the head, in exactly the same
\emph{category} of gap \cref{sec:archB:gap} already names for plain Option B. The honest claim is:
\textbf{bounded and diagnosable distortion, not solved injectivity.} This should be stated exactly this
plainly wherever this pipeline is invoked elsewhere in the note, the same way \cref{sec:archB:gap} states
its own gap without softening it.

\subsection{Gotchas}
\label{sec:pseudoinj:gotchas}
\begin{itemize}[leftmargin=1.5em]
\item \textbf{KL vs.\ injective/bi-Lipschitz tension, and the limits of what can be said about it.} The
      KL term in \cref{eq:pseudoinj-loss} rewards posterior compression toward the prior; the injective and
      bi-Lipschitz terms reward keeping distinct inputs' encodings apart. These pull in opposite
      directions. The most that can honestly be said at present: Zhan et al.\ state that BLAE's losses can
      be added to a VAE, and presumably tested this on a toy example, but \textbf{no such implementation
      was released} (\cref{sec:pseudoinj:blae}), so this note's combination of a fresh
      \texttt{AutoencoderKL} head with BLAE's losses is unverified engineering, not a reproduction of a
      published result. $\lambda_{\mathrm{kl}}$ against $\lambda_{\mathrm{inj}},\lambda_{\mathrm{biLip}}$
      needs to be tuned empirically, with no prior result to anchor the starting point.
\item \textbf{Extreme compression ratio.} \texttt{AutoencoderKL}'s stock configurations target hundreds 
      to thousands of latent channels, not $\sim\!10$; squeezing this far raises real underfitting/
      posterior-collapse risk. \Cref{sec:pseudoinj:arch} above hand-tunes depth and channel width down 
      for this reason, but the $4\times4\times4=64$ scalars $\to20$-scalar bottleneck introduced there 
      is itself a 
      further, unusually aggressive compression step (pooling/flattening plus one linear layer) with no 
      precedent in the stock class to calibrate against.
\item \textbf{Fresh initialization, no pretraining head start.} Unlike IN-VAE, the head is trained from
      scratch on $\sim\!18\mathrm{k}$ TI images alone.
\item \textbf{BLAE losses must be computed on $\hat\mu$, never on the sampled $\hat z$} -- this needs to 
      be explicit in the loss code, not left implicit, since it is easy to accidentally apply it after 
      the reparameterization step.
\item \textbf{No VAE implementation exists in the BLAE repository at all}; combining a tested
      \texttt{AutoencoderKL} \cite{aekl}with a from-scratch reimplementation of BLAE's loss terms is 
      novel engineering that must be built and validated, not assumed to inherit the paper's toy-scale 
      results.
\item \textbf{Spatial, not pooled, embeddings, and a non-standard bottleneck.} IN-VAE's output is a
      $32\times16\times16$ tensor, unlike \nameref{sec:archB}'s flat pooled $\sim\!3072$-dim vector; the
      head must be conv-based to consume this properly, not a bare MLP as in \texttt{app:code}'s
      \texttt{encoder\_mu}/\texttt{encoder\_logvar}. Because stock \texttt{AutoencoderKL} \cite{aekl} 
      has no
      built-in path from a spatial bottleneck to a flat per-image vector, the
      \texttt{quant\_conv}/\texttt{post\_quant\_conv} replacement in \cref{sec:pseudoinj:arch} is
      hand-written and untested code, not a configuration of an existing, validated path.
\item \textbf{Resolution mismatch.} IN-VAE candidates (including the chosen one) are pretrained at 
      $256$px; this note's pipeline currently assumes $224$px throughout (\cref{sec:archB}) and needs 
      updating if this backbone is adopted.
\item \textbf{Screened candidates.} E2E-VAE and MAETok/STELLAR carry alignment or semantic-shaping losses
      that risk discarding style/location signal, the same failure mode \cref{sec:pretrain} already flags
      for supervised ImageNet pretraining; neither is used here, and any candidate -- including the chosen
      IN-VAE -- still needs to clear \cref{app:probe}'s linear-probe screen before being trusted.
\item \textbf{The geodesic-distance table (\texttt{GeoD}) cost, now one stage removed from pixels.} Building
      \texttt{GeoD} on the cached IN-VAE embeddings ($N\!\approx\!18\mathrm{k}$, still $\sim\!8192$-dim)
      is a one-time precompute, run once per choice of backbone/embedding (not per epoch, not per
      minibatch), with two distinct costs:
      \begin{itemize}[leftmargin=1.3em]
      \item \emph{Pairwise-distance step}, $O(N^2d)$. Naive broadcasting would try to materialize an
            $N\times N\times d$ tensor ($\sim\!10$TB at $N=18\mathrm{k}$, $d=8192$) -- infeasible. Use the
            Gram-matrix identity $\|x-y\|^2=\|x\|^2+\|y\|^2-2x^\top y$ (one matmul, $O(Nd+N^2)$ peak memory)
            instead; confirm \texttt{sklearn}'s \texttt{kneighbors\_graph} is honoring its
            \texttt{working\_memory} row-chunking rather than being overridden. Further remedies if still
            too costly: reduce $d$ via PCA/random projection before this step ($8192\to256$,
            Johnson-Lindenstrauss-justified, $\sim\!32\times$ cheaper); approximate kNN (FAISS/HNSW,
            GPU-capable, never materializes a dense pairwise matrix); or a landmark-Isomap variant (exact
            geodesics from a landmark subset only, e.g.\ $2\mathrm{k}$ of $18\mathrm{k}$, nearest-landmark
            interpolation for the rest, turning $O(N^2)$ into $O(N_LN)$).
      \item \emph{All-pairs shortest-path step}, $O(N^2\log N)$ (Dijkstra/Johnson's algorithm over the
            resulting sparse $k$-NN graph). Remedies: shard per-source Dijkstra runs across CPU processes
            (embarrassingly parallel); reuse the landmark trick above to cut the number of source runs from
            $N$ to $N_L$; confirm \texttt{scipy} is selecting Johnson's algorithm (\texttt{method='J'}) for
            the sparse case rather than defaulting to a slower per-node routine.
      \end{itemize}
\item \textbf{Memory.} The dense $N\times N$ \texttt{GeoD} table ($\sim\!1.3$GB at $N=18\mathrm{k}$)
      belongs on CPU (plain array, or \texttt{numpy.memmap} if it outgrows RAM) and should never be copied
      to GPU in full; slicing a small $B\times B$ submatrix per minibatch on CPU before moving only that
      slice to GPU (the pattern already used for \texttt{GeoD} indexing) is correct and should be kept.
      The real risk is upstream, in the pairwise-distance step: doing that naively on GPU can exhaust GPU
      memory before \texttt{GeoD} is even built; the same Gram-matrix trick, or a GPU-batched approximate-
      kNN library, avoids this. If \texttt{GeoD} itself is still tight on a given host, storing it as
      \texttt{float16} (halving it to $\sim\!650$MB at $N=18\mathrm{k}$) is an acceptable fallback, since
      geodesic distances are a regularization target, not a quantity needing \texttt{float32} precision.
\item \textbf{Still not exact injectivity.} As stated in \cref{sec:pseudoinj:principled}, this pipeline
      bounds and quantifies the gap from \cref{thm:ng}'s exact-diffeomorphism hypothesis; it does not
      close it. Any downstream use of this note's identifiability results should carry this caveat forward
      explicitly.
\end{itemize}

% ---------------------------------------------------------------------
% Bibliography entries this section assumes but that are not yet in biblio.tex:
%   zhan2026blae   -- Zhan et al., "BLAE" / bi-Lipschitz autoencoder, arXiv:2604.06701
%   leng2025repae  -- Leng et al., REPA-E (or the specific IN-VAE checkpoint's source paper)
%   (optionally: a citation for SD-VAE/latent-diffusion VAEs and the original
%    VAE paper (Kingma and Welling), if you want them cited by name rather
%    than referred to descriptively as above)
% Also not yet defined in this note: \label{app:code:lossterms} referenced in the
% footnote on the moral-loss term -- point it at the right subsection of app_ng_reference_impl.tex
% (currently that file has no such label).